% TOP_PROBLEM: {"id": 1957, "title": "Erdős Problem #148", "category_id": 1, "category": {"id": 1, "name": "number_theory", "display_name": "Number Theory", "description": "Properties of integers, prime numbers, Diophantine equations.", "slug": "number-theory", "order_index": 1, "created_at": "2026-07-31T15:26:25.670Z"}, "status": "open"}
% FIRST_POSTED: null
% ENTRY_KIND: statement_only
% Imported from: unsolvedmath_additional_500.tex; UnsolvedMath ID 1957
% !TeX program = lualatex
% Compile twice: lualatex 1957.tex
% Self-contained: no external bibliography, images, or \input files.
% Numeric section labels and visible numbers are original UnsolvedMath IDs.
\documentclass[11pt,a4paper]{article}
\usepackage[margin=24mm,headheight=14pt]{geometry}
\usepackage{fontspec}
\setmainfont{Latin Modern Roman}
\setsansfont{Latin Modern Sans}
\setmonofont{Latin Modern Mono}
\usepackage{amsmath,amssymb,mathtools}
\usepackage{unicode-math}
\setmathfont{Latin Modern Math}
\usepackage{microtype}
\usepackage{enumitem,xurl,needspace}
\usepackage[dvipsnames]{xcolor}
\usepackage{hyperref}
\hypersetup{unicode=true,colorlinks=true,linkcolor=MidnightBlue,urlcolor=MidnightBlue,
 pdftitle={UnsolvedMath Problem 1957},pdfauthor={Source-annotated compilation},bookmarksnumbered=true}
\usepackage{bookmark,tocloft,titlesec,fancyhdr}
\setlength{\cftsecnumwidth}{4.2em}
\cftsetpnumwidth{2.5em}
\cftsetrmarg{4em}
\titleformat{\section}{\Large\bfseries\raggedright}{\thesection}{0.7em}{}
\pagestyle{fancy}\fancyhf{}
\fancyhead[L]{\small\sffamily Additional UnsolvedMath statements}
\fancyhead[R]{\small Original catalogue IDs}
\fancyfoot[C]{\thepage}
\setlength{\parindent}{0pt}
\setlength{\parskip}{5pt plus 1pt}
\setlength{\emergencystretch}{3em}
\setcounter{tocdepth}{1}\setcounter{secnumdepth}{1}
\setlist[itemize]{leftmargin=3em,itemsep=4pt,topsep=4pt,parsep=0pt}
\allowdisplaybreaks
\newcommand{\N}{\mathbb{N}}
\newcommand{\Z}{\mathbb{Z}}
\newcommand{\Q}{\mathbb{Q}}
\newcommand{\R}{\mathbb{R}}
\newcommand{\C}{\mathbb{C}}
\newcommand{\F}{\mathbb{F}}
\newcommand{\A}{\mathbb{A}}
\newcommand{\PP}{\mathbb{P}}
\newcommand{\E}{\mathbb{E}}
\newcommand{\eps}{\varepsilon}
\newcommand{\dd}{\,\mathrm d}
\newcommand{\sref}[2]{\hyperlink{source:#1:#2}{[S#2]}}
\newif\ifshowcatalogue
\showcataloguetrue
\newcommand{\cataloguescope}[1]{\ifshowcatalogue
 \paragraph{Statement recorded in the supplied source.}
 #1\fi}
\begin{document}
% UnsolvedMath ID 1957; source catalogue code EP-148

\setcounter{section}{1956}

\section{Counting Egyptian-Fraction Expansions of One}\label{1957}

{\small\sffamily UnsolvedMath ID 1957 · Number Theory}

\subsection{Definitions and mathematical statement}

\paragraph{Shared notation.}Unless a section says otherwise, $\N=\{1,2,\ldots\}$,
$\Z$ is the ring of integers, $\Q,\R,\C$ are the rational, real, and complex
fields, $[N]=\{1,\ldots,\lfloor N\rfloor\}$, and $\log$ is the natural
logarithm. For real functions of a parameter tending to infinity,
$f\ll g$ and $f=O(g)$ mean $|f|\leq Cg$ eventually for some fixed $C$;
$f\gg g$ reverses this comparison; $f=o(g)$ means $f/g\to0$;
$f\sim g$ means $f/g\to1$; and $f\asymp g$ means both order bounds.
A subscript on $O$ or $\ll$ specifies allowed constant dependence.
The expression $n^{o(1)}$ in an upper bound means growth smaller than
$n^\epsilon$ for each fixed $\epsilon>0$. Local definitions govern any
exception, finite-field convention, or use of the same symbol for another
object. An asymptotic claim is not automatically an effective algorithm.



A solution is an increasing tuple of positive integers, so permutations of the same denominators are not counted again. All summands are used exactly once. The task concerns the growth of this finite number of solutions as the length $k$ increases. \sref{1957}{1} \sref{1957}{2}

\paragraph{Statement recorded in the supplied source.}
Let $F(k)$ be the number of solutions to $  1= \frac{1}{n_1}+\cdots+\frac{1}{n_k}, $ where $1\leq n_1<\cdots<n_k$ are distinct integers. Find good estimates for $F(k)$. \sref{1957}{1}

\subsection{Short English statement}

How many increasing lists of a prescribed number of distinct positive denominators have reciprocal sum exactly one?

\subsection{Sources}

{\small

\begin{itemize}

\item[\hypertarget{source:1957:1}{S1}] ULAM.AI, \emph{UnsolvedMath}, supplied \texttt{problems.json}, record 1957 (EP-148), statement and background. The embedded literature-review date is 2026-08-17; that is the archive’s review date, not a fresh certification. Dataset landing page: \url{https://huggingface.co/datasets/ulamai/UnsolvedMath}.

\item[\hypertarget{source:1957:2}{S2}] T. F. Bloom, \emph{Erdős Problems}, Problem 148, \url{https://www.erdosproblems.com/148}. Reference imported from the supplied record; the live problem page was not independently checked for this compilation.

\item[\hypertarget{source:1957:3}{S3}] S. V. Konyagin, Double Exponential Lower Bound for the Number of Representations of Unity by Egyptian Fractions, Mathematical Notes 95 (2014), 277-281, \nolinkurl{doi:10.1134/S0001434614010295.} \url{https://doi.org/10.1134/S0001434614010295}. Bibliographic information imported from the supplied record.

\item[\hypertarget{source:1957:4}{S4}] Christian Elsholtz, Egyptian Fractions with odd denominators, arXiv:1606.02117 (2016). \url{https://arxiv.org/abs/1606.02117}. Bibliographic information imported from the supplied record.

\item[\hypertarget{source:1957:5}{S5}] Elsholtz, Christian and Planitzer, Stefan, Sums of four and more unit fractions and approximate parametrizations. Bull. Lond. Math. Soc. (2021), 695-709. Bibliographic information imported from the supplied record.

\end{itemize}

}

\label{end:1957}
\end{document}
